\subsection{Targets and the degeneracy hypothesis}
\label{subsec:targets-and-the-degeneracy-hypothesis}

We consider a ten-parameter description of a compact binary coalescence signal injected into two-detector noise: chirp mass $\mathcal{M}$, mass ratio $q$, inclination $\iota$, coalescence phase $\varphi_c$, polarization angle $\psi$, declination, right ascension, injection time, merger time within the analysis window, and network signal-to-noise ratio.
Of these, seven are used as supervised regression targets (mass ratio and injection time are excluded; \S\ref{sec:phicpsi:dataset}).
The scientific targets are $\varphi_c \in [0, 2\pi)$, with period $2\pi$, and $\psi \in [0, \pi)$, with period $\pi$.
Inclination, chirp mass, merger time, SNR, and sky position are retained as \emph{control heads}: parameters trained through the identical pipeline whose success or failure calibrates what the network can extract from strain when information is present.

The \textbf{degeneracy hypothesis} under test is: \emph{$\varphi_c$ and $\psi$ carry no strain-only recoverable signal for this population, in this architecture and loss family} --- i.e., a point-estimating network can do no better than the optimal constant (or constant-per-mode) prediction, because for most of the population the likelihood constrains only a combination of the two angles, and the remaining direction is unconstrained.

The alternative hypothesis motivating the experimental design is that even if $\varphi_c$ and $\psi$ are individually unrecoverable, their \textbf{sum and difference combinations} may be well-conditioned: face-on systems constrain $\varphi_c \mp 2\psi$ (sign by $\cos\iota$), so a network parameterized in combination space, with a curriculum that weights each combination according to how well the current inclination constrains it, might learn the recoverable structure that a naive per-angle parameterization misses.

One scope note applies to the whole paper and is stated here rather than left implicit.
Conditioning on inclination $\iota$ --- tested as future work in \S\ref{sec:phicpsi:future} --- is treated throughout as a \emph{training-paradigm design choice}, motivated by the analytic structure of \S\ref{sec:phicpsi:prereq} showing that the well-constrained combination is recoverable given $\iota$, not as a forced move necessitated by a proof that $\iota$ itself is unrecoverable.
No model in this paper takes $\iota$, or any function of it, as an input at any stage; $\iota$ appears here only as a control \emph{output} head trained on raw strain, exactly like $\varphi_c$ and $\psi$ (\S\ref{sec:phicpsi:circular}).
The design that would feed true $(\sin\iota, \cos\iota)$ as an input channel belongs to the successor experiment scoped in \S\ref{sec:phicpsi:future}(i), not to any run reported here.
The inclination head's own failure (\S\ref{sec:phicpsi:pathology2}) is unresolved, not diagnosed as fundamental: it could reflect a genuine unlearnability of $\iota$ from strain under this architecture, or a fixable loss-path choice --- \S\ref{sec:phicpsi:pathology2} traces it to a Huber loss on a raw two-vector, unlike the successful von Mises--Fisher parameterization used for sky position.
Whether $\iota$ can be predicted by an independent pipeline is therefore a separate, open question that this paper does not resolve and does not need to resolve to motivate testing $\iota$-conditioning as a training paradigm; \S\ref{sec:phicpsi:future} excludes inference-time $\iota$ estimation from scope for exactly this reason, not because $\iota$ is assumed to be fundamentally unrecoverable.

\subsection{Circular regression framework}
\label{sec:phicpsi:circular}

Periodic targets are regressed on the unit circle~\citep{mardia2000circular}.
An angle $\theta$ with period $T$ is encoded as the two-vector
\begin{equation}
    \mathbf{u}(\theta) = \bigl(\sin(2\pi\theta/T),\ \cos(2\pi\theta/T)\bigr),
\end{equation}
so the $\psi$ head, with $T=\pi$, physically encodes $2\psi$.
Each periodic head outputs a raw two-vector $\mathbf{v} \in \mathbb{R}^2$, projected onto the unit circle by a normalization layer, \texttt{normalize\_unit},
\begin{equation}
    \hat{\mathbf{u}} = \frac{\mathbf{v}}{\max(\lVert\mathbf{v}\rVert,\ \varepsilon)},
    \qquad \varepsilon = 10^{-8},
\end{equation}
and trained with the cosine (circular) loss
\begin{equation}
    \mathcal{L}_{\mathrm{circ}}
    = 1 - \hat{\mathbf{u}}_{\mathrm{pred}} \cdot \mathbf{u}_{\mathrm{true}}
    = 1 - \cos\Delta\theta .
    \label{eq:circloss}
\end{equation}
For uniformly distributed targets and any fixed prediction, the expected circular loss is exactly 1; the expected wrap-aware angular mean absolute error (ang\_MAE) is $T/4$ --- that is, $\pi/2 \approx {}$\SI{1.5708}{\radian} for $\varphi_c$ and $\iota$, and $\pi/4 \approx {}$\SI{0.7854}{\radian} for $\psi$.
These are the null baselines against which every result in this paper is judged.

Two further diagnostic statistics recur throughout.
The \textbf{circular resultant} of a set of predicted angles, $r_{\mathrm{circ}} = \lVert \operatorname{mean}(\hat{\mathbf{u}}) \rVert$, measures mode collapse: $r_{\mathrm{circ}} \to 1$ means all predictions concentrate at one angle (the empirical baseline for our uniform targets is $r_{\mathrm{circ}} \approx 0.005$).
The \textbf{std\_ratio} of a head is the ratio of the standard deviation of its raw outputs to that of the encoded targets, accumulated over an epoch; because the targets lie on the unit circle, std\_ratio is a direct monitor of raw-magnitude health, with values far from unity signalling the pathology analyzed in \S\ref{sec:phicpsi:pathology2}.
We adopt $[0.5, 2.0]$ as the healthy band.

Non-periodic heads use a Huber loss~\citep{huber1964robust} ($\delta=1$) on transformed targets (log--$z$-scored chirp mass, $z$-scored SNR, sigmoid-bounded merger time), and the sky-position head uses a closed-form von Mises--Fisher negative log-likelihood on the unit sphere~\citep{fisher1953dispersion}.
One head cuts across this taxonomy: the inclination head keeps the periodic two-vector encoding of its target, but its loss is the Huber loss applied directly to the raw output vector --- no \texttt{normalize\_unit} projection and no circular loss anywhere in its path (\S\ref{sec:phicpsi:pathology2} documents the code-level trace and its consequences for $\iota$'s role as a control).
All heads share a single trunk and are combined by learned homoscedastic uncertainty weighting~\citep{kendall2018multi}: each head $h$ carries a trainable log-variance $s_h$, clamped to $|s_h| \le 3$, contributing $e^{-s_h}\mathcal{L}_h + s_h$ to the total loss.